\begin{abstract}
  Constraints allow developers to specify desired properties of
  systems in a number of domains, and have those properties be
  maintained automatically.  This results in compact, declarative
  code, avoiding scattered code to check and imperatively re-satisfy
  invariants.  Despite these advantages, constraint programming is not
  yet widespread, with imperative programming still the norm.

  There is a long history of research on integrating
  constraint programming with the imperative
  paradigm. However, this integration typically does not unify the
  constructs for encapsulation and abstraction from both
  paradigms. This impedes re-use of modules, as client code written in
  one paradigm can only use modules written to support that
  paradigm. Modules require redundant definitions if they are to be
  used in both paradigms.

  We present a language -- \langname~-- that unifies the constructs
  for encapsulation and abstraction
  by using only object-oriented method definitions
  for both declarative and imperative code. Our prototype --
  \toolname~-- is an extension to Ruby, and continues to support
  Ruby's object-oriented semantics.  It allows
  programmers to add constraints to existing Ruby programs in
  incremental steps by placing them on the results of normal
  object-oriented message sends.  It is implemented by modifying
  a state-of-the-art Ruby virtual machine.  The performance of standard
  object-oriented code without constraints is only modestly impacted,
  with typically less than 10\% overhead compared with the unmodified
  virtual machine. Furthermore, our architecture for adding multiple
  constraint solvers allows \langname to deal with constraints in a
  variety of domains.

  We argue that our approach provides a useful step toward making
  constraint solving a generic tool for object-oriented programmers.
  We also provide example applications, written in our Ruby-based
  implementation, which use constraints in a variety of application
  domains, including interactive graphics, circuit simulations, data
  streaming with both hard and soft constraints on performance, and
  configuration file management.

  % We present an approach and extension to an object-oriented language
  % to solve constraints on the results of message sends. These
  % constraints are predicates that are continuously maintained by the
  % system.

  % Constraints are written as anonymous functions that contain
  % constraint expressions. To maintain a constraint, our approach
  % chooses a constraint solver based on the return types of the
  % messages sent in the dynamic extent of the expression execution.
  % Solvers can be flexibly bound to types such as reals or strings, but
  % also complex objects such as files or methods. Furthermore,
  % constraint solvers are implemented in application code, so
  % programmers can add domain specific solvers for their use cases.

  % Our system solves constraints by modifying objects that determine
  % the results of messages sent during the constraint expression
  % execution. Whenever imperative code changes these objects, the
  % appropriate solvers are triggered to maintain the constraints.

  % Our approach is superior to existing approaches because it relies on
  % object-oriented message sends instead of paths between objects or
  % special constraint method definitions. This property ensures both
  % encapsulation and supports the uniform access principle. This, in
  % turn, allows programmers to change the implementation of objects
  % without breaking constraint code, as well as applying groups of
  % constraints to different objects that implement the same
  % interface. Since programmers can extend the available solvers, they
  % can choose the solver that gives the best results or provides the
  % best performance to their specific application.

  % We argue that using our approach, constraint solving becomes a
  % generic tool that is not limited to specific domains. Using a
  % Ruby-based implementation, we show example applications that apply
  % constraints to interesting new domains such as performance
  % assurances and configuration file handling.

\end{abstract}

%%% Local Variables: 
%%% mode: latex
%%% TeX-master: "babelsberg_jot"
%%% End: 
